\chapter{Wstęp}
\label{ch:wstep}

Współczesne systemy pomiarowe, układy automatyki przemysłowej oraz przekształtniki energoelektroniczne wymagają niezawodnego i~precyzyjnego przetwarzania sygnałów analogowych na postać cyfrową[cite: 1, 11]. W~wielu aplikacjach, takich jak pomiar prądów czy napięć w~obwodach mocy, kluczowym parametrem nie jest wyłącznie maksymalna szybkość akwizycji, lecz wysoka rozdzielczość pomiarowa oraz odporność na zakłócenia[cite: 1, 11]. W~tego typu zastosowaniach dużą popularnością cieszą się przetworniki analogowo-cyfrowe (ADC -- ang. \textit{Analog-to-Digital Converter}) wykorzystujące modulację Sigma-Delta ($\Sigma\Delta$)[cite: 1, 11].

Na rynku dostępnych jest wiele zintegrowanych modulatorów $\Sigma\Delta$, takich jak układ AMC1306 firmy Texas Instruments[cite: 18, 23, 30] czy AD7402 firmy Analog Devices[cite: 3, 10]. Oferują one bardzo dobre parametry przetwarzania oraz często wbudowaną barierę izolacji galwanicznej, jednak ich cena jednostkowa bywa wysoka, a~wewnętrzna struktura układu pozostaje zamknięta dla projektanta[cite: 11, 28, 30]. Jednocześnie na rynku podzespołów elektronicznych zauważalna jest luka w~segmencie tańszych układów pomiarowych o~średniej i~wysokiej precyzji, które nie wymagają zaawansowanych funkcji dodatkowych[cite: 28, 30]. Naturalną alternatywą staje się próba wypełnienia tej luki poprzez budowę modulatora $\Sigma\Delta$ z~powszechnie dostępnych, standardowych elementów dyskretnych -- wzmacniaczy operacyjnych, komparatorów oraz pojedynczych bramek i~przerzutników cyfrowych[cite: 7, 25, 30]. Takie podejście pozwala na swobodny dobór komponentów pod kątem kosztowym i~technicznym, jednak wiąże się z~koniecznością uwzględnienia fizycznych ograniczeń poszczególnych elementów, które w~układach dyskretnych są znacznie bardziej widoczne niż wewnątrz dedykowanych struktur krzemowych[cite: 24, 30].

%%%%%%%%%%%%%%%%%%%%%%%%%%%%%%%%%%%%%%%%%%%%%%%%
\section{Zasada działania ciągłych modulatorów Sigma-Delta}
\label{sec:zasada_dzialania}

Klasyczne przetworniki analogowo-cyfrowe (np. o~architekturze sukcesywnej aproksymacji -- SAR) próbkują sygnał wejściowy z~częstotliwością zbliżoną do częstotliwości Nyquista i~bezpośrednio przypisują każdej próbce wielobitowe słowo cyfrowe za pomocą precyzyjnej drabinki napięć odniesienia. Przetworniki $\Sigma\Delta$ realizują to zadanie w~odmienny sposób: zamiast rozbudowanego, wielobitowego kwantyzatora analogowego wykorzystują prosty kwantyzator jednobitowy (komparator) wspierany przez szybką logikę cyfrową[cite: 1, 6, 16]. Wysoką rozdzielczość końcową uzyskuje się dzięki połączeniu dwóch mechanizmów: nadpróbkowania (ang. \textit{oversampling}) oraz kształtowania widma szumu kwantyzacji (ang. \textit{noise shaping})[cite: 1, 16, 30].

Nadpróbkowanie polega na próbkowaniu sygnału wejściowego z~częstotliwością $f_s$ wielokrotnie wyższą od maksymalnej częstotliwości użytecznej mierzonego sygnału[cite: 16]. Sprawia to, że całkowita moc szumu kwantyzacji rozkłada się na bardzo szerokie pasmo częstotliwości, dzięki czemu tylko niewielka jej część przypada na wąskie pasmo sygnału użytecznego[cite: 16]. Drugi mechanizm -- kształtowanie szumu -- wynika z~zamknięcia kwantyzatora w~ujemnej pętli sprzężenia zwrotnego zawierającej filtr pętli (układ całkujący)[cite: 1, 2]. Dla niskich częstotliwości (pasmo sygnału) układ zachowuje się jak wzmacniacz o~stałym wzmocnieniu, natomiast dla błędu kwantyzacji pętla stanowi filtr górnoprzepustowy, który wypycha szum kwantyzacji w~obszar wysokich częstotliwości, poza pasmo użyteczne[cite: 1, 2, 16].

Ze względu na sposób realizacji filtru pętli modulatory $\Sigma\Delta$ dzieli się na układy z~czasem dyskretnym (DT -- ang. \textit{Discrete-Time}), realizowane najczęściej w~technice przełączanych pojemności, oraz układy z~czasem ciągłym (CT -- ang. \textit{Continuous-Time})[cite: 1, 2, 4]. W~niniejszej pracy analizie poddano architekturę z~czasem ciągłym (CT), której ogólny schemat blokowy przedstawiono na rysunku~\ref{rys:ct_sdm_ogolny}[cite: 2, 6, 30].

\begin{figure}[!htbp]
	\centering
	% Miejsce na schemat blokowy ciągłego modulatora Sigma-Delta (np. wyeksportowany do PDF/PNG)
	% \includegraphics[width=0.75\linewidth]{ct_sdm_schemat.pdf}
	\fbox{\parbox{0.75\linewidth}{\centering \vspace{1.5cm} Miejsce na rysunek: Ogólny schemat blokowy ciągłego modulatora $\Sigma\Delta$ (integrator analogowy, komparator, przerzutnik D oraz pętla sprzężenia zwrotnego) \vspace{1.5cm}}}
	\caption{Ogólny schemat blokowy modulatora Sigma-Delta z~czasem ciągłym (CT)[cite: 2, 6]}
	\label{rys:ct_sdm_ogolny}
\end{figure}

W~modulatorze ciągłym (rysunek~\ref{rys:ct_sdm_ogolny}) analogowy sygnał wejściowy nie jest próbkowany na samym wejściu układu, lecz trafia bezpośrednio na analogowy układ całkujący (integrator zbudowany na wzmacniaczu operacyjnym z~elementami $R$ oraz $C$)[cite: 2, 6, 25]. W~węźle sumującym od sygnału wejściowego odejmowany jest sygnał sprzężenia zwrotnego[cite: 2, 6]. Przecałkowana różnica sygnałów (błąd) podawana jest na komparator, który pełni rolę jednobitowego przetwornika analogowo-cyfrowego, a~następnie na przerzutnik typu D taktowany zegarem o~częstotliwości próbkowania $f_s$[cite: 6, 25]. Dopiero na przerzutniku następuje dyskretyzacja sygnału w~czasie[cite: 2, 6]. Wyjście przerzutnika steruje pętlą sprzężenia zwrotnego (pełniąc rolę jednobitowego przetwornika cyfrowo-analogowego) i~jednocześnie stanowi wyjście całego modulatora[cite: 2, 6, 25]. Istotną zaletą architektury ciągłej jest to, że znajdujący się na wejściu układ całkujący działa jak naturalny filtr antyaliasingowy, tłumiący zakłócenia o~wysokich częstotliwościach jeszcze przed procesem próbkowania[cite: 3, 6].

Zwiększenie liczby stopni całkowania w~pętli sprzężenia zwrotnego (np. przejście z~modulatora 1.~rzędu na modulator 2.~rzędu z~dwoma kaskadowo połączonymi integratorami) pozwala na bardziej strome kształtowanie widma szumu kwantyzacji i~uzyskanie wyższej rozdzielczości przy tym samym współczynniku nadpróbkowania, kosztem większej złożoności układowej i~mniejszego zapasu stabilności[cite: 3, 6, 16]. Surowym sygnałem wyjściowym każdego z~omawianych modulatorów jest wysokoczęstotliwościowy strumień pojedynczych bitów (ang. \textit{bitstream}), w~którym informacja o~chwilowej wartości napięcia wejściowego zakodowana jest w~gęstości występowania stanów wysokich[cite: 6, 16]. Aby odzyskać z~niego wielobitową wartość pomiarową oraz odfiltrować przesunięty w~górę pasma szum kwantyzacji, strumień ten poddaje się filtracji dolnoprzepustowej i~decymacji w~układzie cyfrowym (np. w~strukturze FPGA lub mikrokontrolerze) za pomocą filtrów typu Sinc (CIC -- ang. \textit{Cascaded Integrator-Comb})[cite: 16, 19, 30].

%%%%%%%%%%%%%%%%%%%%%%%%%%%%%%%%%%%%%%%%%%%%%%%%
\section{Cel, założenia i~zakres pracy}
\label{sec:cel_pracy}

Głównym celem niniejszej pracy jest symulacyjna oraz sprzętowa ocena efektywności i~zasadności budowy ciągłych modulatorów Sigma-Delta z~elementów dyskretnych[cite: 30]. Projekt dąży do udzielenia odpowiedzi na pytanie, jaką efektywną rozdzielczość (wyrażoną w~bitach) można uzyskać, budując modulator z~niedrogich, ogólnodostępnych podzespołów elektronicznych, oraz które parametry rzeczywistych elementów stanowią główne ograniczenie tego rozwiązania w~zestawieniu z~komercyjnym układem scalonym Texas Instruments AMC1306[cite: 24, 30].

Aby zrealizować postawiony cel główny, przyjęto następujący zakres prac projektowych i~badawczych[cite: 30]:
\begin{itemize}
	\item opracowanie w~środowisku MATLAB/Simulink (z~wykorzystaniem pakietu PLECS Blockset) czterech modeli symulacyjnych: idealnego modulatora 2.~rzędu, modulatora 1.~rzędu zbudowanego z~tanich elementów dyskretnych oraz dwóch wariantów modulatora 2.~rzędu (wykorzystujących odpowiednio tanie oraz droższe, szybsze komponenty analogowe)[cite: 24, 30],
	\item zamodelowanie nieidealności fizycznych podzespołów: skończonego wzmocnienia w~otwartej pętli ($A_{\text{OL}}$), skończonego pola wzmocnienia (GBW), szybkości narastania napięcia (ang. \textit{slew rate}), napięcia niezrównoważenia (offsetu) i~nasycenia wzmacniaczy operacyjnych, a~także histerezy oraz czasów propagacji komparatora i~przerzutnika[cite: 24, 30],
	\item przeprowadzenie wielowymiarowej analizy parametrycznej określającej wpływ poszczególnych niedoskonałości układowych oraz stałych czasowych integratorów na wypadkowy stosunek sygnału do szumu i~zniekształceń (SINAD) oraz efektywną liczbę bitów (ENOB)[cite: 24, 30],
	\item zaprojektowanie i~wykonanie czterowarstwowego obwodu drukowanego (PCB) integrującego badane warianty modulatorów dyskretnych wraz z~referencyjnym układem scalonym TI AMC1306[cite: 30],
	\item oprogramowanie układu FPGA do równoczesnej akwizycji czterech strumieni bitowych, ich pakietowania za pomocą kolejek FIFO, buforowania w~zewnętrznej pamięci RAM oraz przesyłania przez interfejs szeregowy RS do komputera PC[cite: 30],
	\item przeprowadzenie pomiarów na stanowisku fizycznym, porównanie wyników eksperymentalnych z~symulacjami oraz wykonanie analizy techniczno-ekonomicznej opracowanych układów[cite: 30].
\end{itemize}

Kluczowym założeniem badawczym przyjętym w~części eksperymentalnej jest umieszczenie wszystkich badanych torów (dyskretnych oraz zintegrowanego układu odniesienia TI AMC1306) na wspólnej płytce PCB i~jednoczesne próbkowanie przez nie dokładnie tego samego sygnału wejściowego[cite: 30]. Wyeliminowanie różnic wynikających z~odmiennych warunków zasilania, źródeł sygnału czy zegara taktującego pozwala na bezpośrednie i~miarodajne porównanie jakości konwersji każdego z~układów[cite: 30].

%%%%%%%%%%%%%%%%%%%%%%%%%%%%%%%%%%%%%%%%%%%%%%%%
\section{Struktura pracy}
\label{sec:struktura_pracy}

Praca została podzielona na pięć rozdziałów, których układ odzwierciedla kolejne etapy realizacji projektu inżynierskiego -- od wyznaczenia kryteriów oceny i~badań symulacyjnych, przez budowę sprzętu, aż po weryfikację pomiarową[cite: 30].

W~rozdziale pierwszym zarysowano tło problematyki oraz motywację podjęcia tematu, omówiono ogólną zasadę działania ciągłych modulatorów Sigma-Delta, a~także sformułowano cel i~zakres pracy[cite: 30].

W~rozdziale drugim przedstawiono badania symulacyjne oraz analizę parametryczną zaprojektowanych układów[cite: 30]. Rozdział ten rozpoczyna się od wprowadzenia metodyki oceny jakości sygnałów za pomocą transformaty Fouriera (FFT), definicji parametrów SINAD oraz ENOB, a~także opisu zastosowanego filtru cyfrowego typu Sinc/CIC[cite: 24, 30]. Następnie omówiono strukturę modeli opracowanych w~środowisku Simulink/PLECS, sposób implementacji ograniczeń sprzętowych oraz wyniki analizy pokazującej wagę wpływu poszczególnych parametrów komponentów na efektywność konwersji[cite: 24, 30].

Rozdział trzeci poświęcono projektowi i~realizacji sprzętowej stanowiska badawczego[cite: 30]. Zaprezentowano w~nim schemat blokowy systemu, uzasadniono dobór elementów dla czterech wariantów modulatorów oraz zestawiono koszty budowy układów dyskretnych z~ceną układu referencyjnego[cite: 30]. W~dalszej części rozdziału opisano proces projektowania obwodu drukowanego z~uwzględnieniem zasad prowadzenia sygnałów o~wysokiej częstotliwości, a~także architekturę systemu akwizycji, buforowania i~ekstrakcji danych zaimplementowaną w~układzie FPGA[cite: 30].

W~rozdziale czwartym skupiono się na pomiarach i~weryfikacji działania fizycznego prototypu[cite: 30]. Zarejestrowane strumienie bitowe poddano obróbce za pomocą tych samych algorytmów filtracji i~oceny widmowej, które wykorzystano na etapie symulacji[cite: 30]. Porównano zbieżność wyników pomiarowych z~modelami symulacyjnymi oraz zestawiono parametry układów dyskretnych z~układem scalonym TI AMC1306, oceniając zasadność techniczną i~ekonomiczną obu rozwiązań[cite: 30].

Rozdział piąty stanowi podsumowanie pracy, w~którym zebrano najważniejsze wnioski z~badań symulacyjnych i~eksperymentalnych, omówiono napotkane trudności projektowe oraz wskazano możliwe kierunki dalszego rozwoju zbudowanego stanowiska[cite: 30].